\documentclass[11pt]{article}
\usepackage{mathblatt}
% Lösungen nach mathe-nachhilfe gemeinsam.md „Lösungen“ und pruefung.md § 6.
\newcommand{\kennung}{LXA}
\begin{document}
\pfheftstil
\pfheftkopf[\kennung]{Seite im allgemeinen Dreieck, Sinussatz}{P10 \textperiodcentered{} Lösungen}

\begin{pfloesung}{}
\lz{1.}{$\beta = 110^\circ$, $\gamma = 30^\circ$}{Nebenwinkel: $\beta = 180^\circ - 70^\circ$; Winkelsumme: $\gamma = 180^\circ - 40^\circ - 110^\circ$}{}
\end{pfloesung}
\begin{pfloesung}{}
\lz{2.}{$34^\circ$}{im Dreieck $ABC$ zählt $\beta_2$, nicht $\beta_1$: $180^\circ - 38^\circ - 108^\circ$}{}
\end{pfloesung}
\begin{pfloesung}{}
\lz{3.}{B}{$b$ liegt $\beta$ gegenüber; von $c$ fehlt die Länge}{}
\end{pfloesung}
\begin{pfloesung}{}
\lz{4a)}{$b \approx 6{,}0$ cm}{$b = \dfrac{8 \cdot \sin 45^\circ}{\sin 70^\circ}$}{}
\lz{b)}{$a \approx 5{,}6$ cm}{$\gamma = 68^\circ$; $a = \dfrac{7 \cdot \sin 48^\circ}{\sin 68^\circ}$}{}
\end{pfloesung}
\begin{pfloesung}{}
\lz{5.}{$s \approx 4{,}4$ m}{$\beta = 110^\circ$, $\gamma = 32^\circ$; $s = \dfrac{2{,}5 \cdot \sin 110^\circ}{\sin 32^\circ}$}{}
\end{pfloesung}
\begin{pfloesung}{}
\lz{6.}{$y \approx 180{,}1$ cm}{dritter Winkel $34^\circ$; $y = \dfrac{160 \cdot \sin 141^\circ}{\sin 34^\circ}$}{}
\end{pfloesung}
\end{document}
